\documentclass[10pt,twocolumn,letterpaper]{article}
\usepackage{../../templates/singletai-preprint}
\usepackage{booktabs}
\usepackage{longtable}
\usepackage{array}
\usepackage{multirow}
\usepackage{algorithm}
\usepackage{algorithmic}

\title{Singlet AI: Personalized Genomics Intelligence\\
  from Single-Cell Resolution Foundation Models}

\author{%
  Zach DeBruine\textsuperscript{1,*}\\[4pt]
  \textsuperscript{1}Grand Valley State University, Allendale, MI, USA\\
  \textsuperscript{*}Correspondence: \texttt{debruinz@gvsu.edu}
}

\date{\today}

\begin{document}
\maketitle

\begin{abstract}
We present the investment thesis and product strategy for Singlet AI, a company
building the genomics intelligence layer for modern medicine. Our core technical
asset is the world's largest uniformly reprocessed single-cell atlas: 1.5~billion
cells across 19,790 GEO series, 201 species, and 6 modalities---all processed
from raw FASTQ reads through a single pipeline. We propose Singlet PGI
(Personalized Genomics Intelligence), a post-trained frontier AI agent that
translates any patient genome into mechanistic, cell-type-specific explanations of
genetic variants, grounded in interpretable NMF biological programs. The
architecture freezes a sequence-to-function trunk (Borzoi or equivalent), trains a
novel single-cell expression prediction head on paired WGS~+~scRNA-seq data
(TenK10K), projects predictions through multimodal NMF factor models, and
post-trains LLMs to reproduce clinical genetics case reports. We describe the
market opportunity (\$175B precision medicine by 2030), competitive positioning,
technical roadmap, de-risking strategy, and customer segments across clinical
genetics, pharmaceutical R\&D, academic research, and core facility operations.
\end{abstract}

%% ──────────────────────────────────────────────────────────────────
\section{Introduction: Four Convergences}
\label{sec:convergences}

Four independent advances have converged to make personalized genomics intelligence
possible for the first time:

\begin{enumerate}
  \item \textbf{Sequence-to-function models have learned the grammar of genetic
    variation.} A new class of deep neural networks---Borzoi (Calico, Apache~2.0),
    Enformer (DeepMind), AlphaGenome (DeepMind, Nature 2026)---predict molecular
    phenotypes directly from DNA sequence. These models take raw DNA as input and
    predict functional outputs (gene expression, splicing, chromatin accessibility)
    at high resolution. Their convolutional~+ transformer trunks have learned the
    regulatory grammar of the genome. This is fundamentally different from DNA
    language models like Evo~2 (Arc Institute), which learn sequence
    \emph{patterns} rather than sequence-to-function mappings.

  \item \textbf{Paired WGS~+ single-cell data now exists at scale.} TenK10K
    (Cuomo et al., medRxiv 2025) produced 1,925~individuals with matched
    whole-genome sequencing and scRNA-seq: 5~million+ cells across 28 immune cell
    types. They demonstrated that $\sim$50\% of genetic regulatory effects are
    cell-type-specific---invisible to bulk approaches.

  \item \textbf{Single-cell data at atlas scale.} We have uniformly reprocessed all
    19,790 series from NCBI GEO---1.17~million samples, 201~species, 6
    modalities---into the world's largest harmonized single-cell repository, 10$\times$
    the coverage of CZ~CELLxGENE Census.

  \item \textbf{LLMs that reason about biology.} CZI's rbio1 (Istrate et al., 2025)
    proved that frontier LLMs can be post-trained with biological world models as
    reward signals, achieving state-of-the-art perturbation prediction. We extend
    this paradigm with interpretable NMF factors that create reasoning chains
    clinicians can trust.
\end{enumerate}

\subsection{The Rare Disease Beachhead}

We start with rare genetic disorders because they offer the clearest validation
signal. Single variants cause clear phenotypes (high penetrance), an enormous
unmet need exists (25--30M Americans, 5--7~year average diagnostic odyssey),
variants of uncertain significance (VUS) are the core pain point, whole-genome
sequencing adoption is growing rapidly, and PubMed contains thousands of clinical
case studies that serve as natural benchmarks for evaluation.

%% ──────────────────────────────────────────────────────────────────
\section{The Product}
\label{sec:product}

\subsection{Singlet PGI---Personalized Genomics Intelligence}

A post-trained frontier AI agent that takes a patient's genome and returns
mechanistic, cell-type-specific explanations of genetic variants in natural
language a clinician can act on.

\paragraph{For Clinicians.} Given a novel missense variant in BMPR1B with
craniofacial anomalies, PGI returns: affected cell types (cranial neural crest
cells, maxillary mesenchyme), disrupted pathways (BMP-SMAD in osteoblast
differentiation), similar known pathogenic variants (BMPR1A, ACVR1), suggested
clinical workup (3D CT craniofacial, echocardiogram), confidence scores grounded
in 14~supporting datasets and 7~published case reports.

\paragraph{For Pharma.} Given a PCSK9 inhibitor showing unexpected hepatotoxicity
in Phase~II, PGI identifies off-target cholangiocyte ER stress driven by NMF
factor~147 (bile acid transport / ER homeostasis), supported by 3 concordant
perturbation datasets.

\subsection{The Intelligence Layer---Many Verticals, One Engine}

\begin{table}[h]
\centering
\small
\begin{tabular}{@{}llr@{}}
\toprule
\textbf{Vertical} & \textbf{Question} & \textbf{Market} \\
\midrule
Clinical Genetics & What does this variant do? & \$30B \\
Pharma R\&D & Which cell types affected? & \$10B+ \\
DTC Genomics & Cellular-resolution meaning & \$3B \\
Diagnostics & Phenotype from genotype? & \$12B \\
Clinical Decision & Monitoring requirements? & \$15B \\
\bottomrule
\end{tabular}
\caption{Application verticals served by a single intelligence engine.}
\label{tab:verticals}
\end{table}

%% ──────────────────────────────────────────────────────────────────
\section{Technical Architecture}
\label{sec:architecture}

\subsection{The Five-Step Pipeline}

Every prediction follows five steps. The key insight: sequence-to-function models
have learned the grammar of genetic variation, but they operate at bulk
resolution. We add the single-cell layer.

\paragraph{Step 1: DNA $\to$ Single-Cell Transcriptomes.} We freeze a
sequence-to-function trunk (Borzoi, Apache~2.0, fully open weights) and train a
new prediction head that maps DNA embeddings directly to splicing-resolved
single-cell expression profiles (S/U/A layers) across cell types, with
microbiome/virome context (Kraken2). Training data: TenK10K (1,925~donors,
5M+~cells, 28~immune cell types). This is not eQTL mapping---it is a nonlinear
neural network learning full sequence context.

\paragraph{Step 2: Multimodal Projection.} Hierarchical NMF models (trained on
1.5B~cells, 6~modalities, 201~species) project the predicted transcriptome into
chromatin accessibility (ATAC), surface protein abundance (CITE-seq), spatial
tissue organization (Visium), RNA velocity / cell dynamics (S/U/A layers),
perturbation response signatures (Perturb-seq), and microbiome/virome
interactions.

\paragraph{Step 3: Biological Interpretation.} NMF factors map to named
biological programs: which cell types are affected, which tissues, which
developmental time points, how spatial signaling is altered, and what splicing
programs are disrupted.

\paragraph{Step 4: Post-Trained AI Reasoning.} Rich context from Steps~1--3 feeds
frontier LLMs post-trained to reproduce published clinical genetics case reports.
NMF factors serve as soft verifiers (reward signals), grounding the LLM in real
biological data at every reasoning step.

\paragraph{Step 5: Clinical Intelligence Output.} Natural language interface
delivering mechanistic explanations, confidence scores, similar known variants,
suggested clinical workup, and cited evidence. Delivery via Epic App Orchard,
clinical web portal, and API.

\subsection{Why This Architecture Works}

The critical innovation is going directly from DNA sequence to splicing-resolved
single-cell transcriptomes---not through a bulk intermediary---and then using NMF
factor models to reason across every modality relatable to the transcriptome.

\begin{table}[h]
\centering
\small
\begin{tabular}{@{}lll@{}}
\toprule
\textbf{Capability} & \textbf{Required For} & \textbf{Status} \\
\midrule
DNA sequence grammar & Step 1 & Borzoi (open) \\
Paired WGS+scRNA & Step 1 & TenK10K \\
SC atlas (1.5B cells) & Step 2 & \textbf{Our asset} \\
Interpretable NMF & Steps 2--3 & \textbf{Our models} \\
S/U/A layers & Steps 2--3 & \textbf{Unique to us} \\
Biological LLM & Step 4 & rbio1 validates \\
Clinical knowledge & Step 5 & ClinVar, PubMed \\
\bottomrule
\end{tabular}
\caption{Architecture capability matrix.}
\label{tab:capabilities}
\end{table}

\subsection{The TenK10K Bridge}

Blood is the ideal entry point: it is the most commonly collected clinical
specimen, contains 28+ immune cell types reflecting systemic biology, and is
where the largest paired WGS~+ scRNA-seq datasets exist. TenK10K Phase~1
provides 1,925~donors with whole-genome sequencing (not arrays---full rare
variant resolution), 5M+~cells, regulatory signals for 47\% of genes tested
(rare variants) and 83\% of all 21,404 genes (common variants).

\subsection{Five AI Reasoning Innovations}

Our architecture extends the rbio1 paradigm with five innovations:
\begin{enumerate}
  \item \textbf{Genome as input, not perturbation.} The sequence-to-function trunk
    translates DNA variants into functional perturbation vectors.
  \item \textbf{NMF factors as reasoning substrate.} The LLM reasons over
    biological programs (e.g., ``cardiac repolarization'') rather than
    $\sim$20,000 individual genes.
  \item \textbf{NMF factors as soft verifiers.} Dense, grounded biological
    verification at every reasoning step---interpretable unlike scGPT/Geneformer
    embedding distances.
  \item \textbf{Splicing and dynamics as unique signal.} RNA velocity and
    dynamical information predict not just where a cell is but where it is going.
  \item \textbf{Clinical literature as evaluation target.} PubMed case studies
    serve as held-out evaluation for whether the LLM can connect genomic
    perturbations to clinical phenotypes.
\end{enumerate}

%% ──────────────────────────────────────────────────────────────────
\section{Technology Stack}
\label{sec:technology}

\subsection{The Data Asset}

\begin{table}[h]
\centering
\small
\begin{tabular}{@{}lrr@{}}
\toprule
\textbf{Metric} & \textbf{Our Atlas} & \textbf{CELLxGENE} \\
\midrule
GEO series & 19,790 & $\sim$2,080 \\
Total samples & 1,166,887 & --- \\
Estimated cells & $\sim$1.5B & $\sim$149M \\
Species & 201 & $\sim$10 \\
Modalities & 6 & 1 (RNA) \\
S/U/A layers & Yes & No \\
Kraken2 profiling & Yes & No \\
Uniform reprocessing & Yes (FASTQ) & No \\
\bottomrule
\end{tabular}
\caption{Data asset comparison.}
\label{tab:data_comparison}
\end{table}

Unique capabilities include: spliced/unspliced/ambiguous layers for every sample
via simpleaf/alevin-fry (enabling RNA velocity), Kraken2 contamination profiling
per cell, unified gene references (GENCODE~v44 human, vM33 mouse, Ensembl for
all 201~species), 48+ file formats handled, and 99.5\% success rate across
processed series.

\subsection{StreamPress}

Custom compressed sparse format delivering $\sim$10$\times$ compression versus
standard formats, with native GPU streaming and zero-copy PyTorch tensor loading.
The entire 1.5B-cell atlas fits in $\sim$7~TB (StreamPress) instead of 70+~TB.
This drives 93--96\% gross margins.

\subsection{FactorNet---Fast Interpretable NMF}

Out-of-core NMF on GPU and CPU scaling to billions of cells, with hierarchical
factorization, automatic rank determination, and $\sim$100$\times$ faster
performance than transformer foundation models at comparable accuracy. CZI chose
to host our NMF embeddings on CELLxGENE Census alongside Geneformer, scGPT,
TranscriptFormer, and UCE---the only non-transformer model.

\subsection{Multimodal Integration}

The NMF factor space connects all modalities because we train multimodal NMF
models on transcriptomes \emph{in combination with} ATAC, CITE-seq, spatial, and
perturbation data. Any prediction in transcriptome space automatically projects
into every other modality through the shared factor representation.

%% ──────────────────────────────────────────────────────────────────
\section{Market Opportunity}
\label{sec:market}

\subsection{Total Addressable Market}

\begin{table}[h]
\centering
\small
\begin{tabular}{@{}lrrl@{}}
\toprule
\textbf{Market} & \textbf{2024} & \textbf{2030} & \textbf{Role} \\
\midrule
Precision medicine & \$73B & \$175B & Intelligence layer \\
Rare disease & \$8B & \$17B & Beachhead \\
Clinical genomics & \$12B & \$30B & Variant interpret. \\
AI drug discovery & \$1.5B & \$10B+ & Target discovery \\
Single-cell analysis & \$4.9B & \$13.7B & Data infrastructure \\
\bottomrule
\end{tabular}
\caption{Total addressable market by segment.}
\label{tab:tam}
\end{table}

\subsection{Revenue Projections}

\begin{table}[h]
\centering
\small
\begin{tabular}{@{}lrrr@{}}
\toprule
\textbf{Year} & \textbf{Intelligence} & \textbf{Data Infra} & \textbf{Total ARR} \\
\midrule
1 & --- & \$220K & \$220K \\
2 & \$200K & \$930K & \$1.1M \\
3 & \$1.5M & \$3.0M & \$4.5M \\
4 & \$5.0M & \$5.0M & \$10.0M \\
5 & \$12.0M & \$7.0M & \$19.0M \\
\bottomrule
\end{tabular}
\caption{Five-year revenue projection.}
\label{tab:revenue}
\end{table}

%% ──────────────────────────────────────────────────────────────────
\section{De-Risking Strategy}
\label{sec:derisk}

Each milestone creates independent value with multiple de-risking pathways:

\paragraph{De-Risk 1: Single-Cell Data Infrastructure (Tempus Playbook).}
Tempus built a \$6B company by organizing biomedical data at bulk resolution.
We do the same at single-cell resolution---10$\times$ more data than CELLxGENE,
uniformly processed, multimodal.

\paragraph{De-Risk 2: LLM-Annotated Data Products.} Our atlas + LLM annotation
creates uniquely rich metadata: cell type identity, cell cycle stage, donor
demographics, developmental time point, disease state, perturbation conditions,
and differential expression signatures---all standardized across 19,790 datasets.

\paragraph{De-Risk 3: NMF Inference API.} Fast, interpretable cell-type
decomposition as a service at $\sim$100ms per query, powering virtual cell
generation from NMF factors.

\paragraph{De-Risk 4: Intermediate Deliverables.} Multiple products on the path
to clinical PGI: annotated atlas API~(Month~4--6), NMF inference API~(Month~6--8),
cross-modal prediction~(Month~8--12), genome$\to$transcriptome API~(Month~10--14),
pharma perturbation tool~(Month~14--18), and full clinical PGI~(Month~18--30).

%% ──────────────────────────────────────────────────────────────────
\section{Competitive Landscape}
\label{sec:competition}

\subsection{The Core Differentiator}

No existing system goes directly from genome~$\to$ single cell~$\to$ clinical
phenotype. Current clinical genomics AI performs database lookup: if a variant is
known, it returns an actionable mutation; otherwise, 40--60\% of variants end as
VUS. Singlet PGI traverses the full path: genome~$\to$ single-cell head~$\to$
NMF multimodal projection~$\to$ post-trained LLM~$\to$ mechanistic clinical
explanation.

\subsection{Data Infrastructure Competitors}

CZ~CELLxGENE Census (149M cells, 2,080 datasets, free, \$30B+ CZI endowment)
is the primary benchmark but lacks uniform FASTQ reprocessing, S/U/A layers,
200+ species coverage, and multimodal linking. TileDB (\$142M raised) provides
storage infrastructure. 10x~Genomics focuses on instrument sales.

\subsection{Clinical Genomics Competitors}

Tempus AI (\$632M revenue, bulk genomics only), Foundation Medicine (targeted
panels, cancer only), SOPHiA Genetics (bioinformatics pipeline, database lookup),
Fabric Genomics (database lookup). None operate at single-cell resolution.

\subsection{Why CZI Is a Collaborator}

CZI is philanthropic, focused on community collaboration. They have no clinical
product ambitions. Our NMF is already hosted on Census. We could be a data
supplier---our uniformly reprocessed data flowing into CELLxGENE as higher-quality
inputs.

%% ──────────────────────────────────────────────────────────────────
\section{Moat and Defensibility}
\label{sec:moat}

\subsection{Six-Layer Defense}

The technology stack forms a six-layer defense:
Layer~1 (Sequence grammar trunk---commoditized),
Layer~2 (Single-cell prediction head---\textbf{our key innovation}),
Layer~3 (Multimodal atlas---\textbf{data moat}),
Layer~4 (Interpretable NMF---\textbf{our models}),
Layer~5 (Reasoning AI),
Layer~6 (Clinical product).
Nobody else has Layers~2--4 simultaneously.

\subsection{The Data Moat}

$\sim$1.5~billion cells across 201~species, 6~modalities, 48+ formats, all one
pipeline. S/U/A layers for every sample. Kraken2 contamination profiles per cell.
LLM-extracted metadata (20 structured fields per dataset). Multimodal bridges
linking RNA~$\leftrightarrow$ ATAC~$\leftrightarrow$ protein~$\leftrightarrow$
spatial through NMF factors. Replicating this asset requires 600K+ core-hours
and 1--2~years of pipeline development.

%% ──────────────────────────────────────────────────────────────────
\section{Business Model}
\label{sec:business}

\subsection{Pricing}

\begin{table}[h]
\centering
\small
\begin{tabular}{@{}llr@{}}
\toprule
\textbf{Tier} & \textbf{Target} & \textbf{Price} \\
\midrule
Free & Students & \$0 \\
Pro & Researchers & \$99--\$299/mo \\
Team & Labs, biotechs & \$999--\$2,999/mo \\
Enterprise & Pharma, CROs & \$50K--\$200K/yr \\
Clinical PGI & Hospitals & Per-query + sub. \\
\bottomrule
\end{tabular}
\caption{Pricing tiers.}
\label{tab:pricing}
\end{table}

\subsection{Unit Economics}

StreamPress compression drives margin advantage: 96\% gross margin on Pro,
96\% on Team, 93\% on Enterprise.

%% ──────────────────────────────────────────────────────────────────
\section{Technical Roadmap and Budget}
\label{sec:roadmap}

\subsection{The \$50K Foundation}

With $\sim$\$50K of compute, we complete the entire annotated data asset:
RNA-seq reprocessing (\$8,250), ATAC-seq + Multiome (\$5,400), Visium + CITE-seq
(\$3,300), processing contingency (\$4,200), LLM metadata extraction (\$300),
LLM cell-type annotation (\$22,000), and hierarchical NMF training (\$3,000).

\subsection{Phase 1: Data Infrastructure (Months 1--6)}

Complete atlas reprocessing ($\sim$1.4M core-hours, \$21K), LLM annotation
(\$20K--\$50K), GPU model training (\$9K--\$39K), web app development
(\$10K--\$20K). Deliverable: singlet.ai public beta.

\subsection{Phase 2: Scale + Bridge (Months 6--12)}

AWS production infrastructure, single-cell head training on seq-to-function trunk
(\$10K--\$25K), LLM post-training (\$20K--\$50K). Deliverable: revenue from data
infrastructure + genome$\to$clinical reasoning prototype.

\subsection{Phase 3: Clinical PGI (Months 12--24)}

Extend to non-blood tissues, clinical evaluation harness, PGI alpha, clinical
partner pilots, Epic App Orchard, FDA regulatory analysis. Total: \$75K--\$175K.

\subsection{Total Budget Summary}

Total Year~1 (low--high): \$520K--\$1.02M. Compute-only: \$163K--\$374K.

%% ──────────────────────────────────────────────────────────────────
\section{Team and Credibility}
\label{sec:team}

\subsection{Advisors}

Josh Welch (U.~Michigan)---creator of LIGER, the leading NMF-based single-cell
integration method. Peter Laird (Van Andel Institute)---Distinguished
Investigator, led The Cancer Genome Atlas (TCGA).

\subsection{Validation}

$\sim$\$400K in CZI grant funding (Cycle~1 \& 3). NIH Transformative R01 at Van
Andel Institute. NMF hosted on CZ~CELLxGENE alongside Geneformer, scGPT,
TranscriptFormer, UCE. NeurIPS and CZ~Benchmarks presentations.

%% ──────────────────────────────────────────────────────────────────
\section{Go-to-Market}
\label{sec:gtm}

\textbf{Phase 1 (Months 0--6):} Bottom-up academic adoption---free tier, ISMB
and NeurIPS presentations, preprint, CLI + Python client.
\textbf{Phase 2 (Months 6--12):} Enterprise sales---pharma pilots, advisor
introductions, benchmarking service.
\textbf{Phase 3 (Months 12--24):} Clinical entry---academic medical center
partners, rare disease validation, Epic integration.
\textbf{Phase 4 (Month 24+):} Clinical scale---hospital deployments, DTC
partnerships, international expansion.

%% ──────────────────────────────────────────────────────────────────
\section{The Ask}
\label{sec:ask}

\textbf{\$500K--\$1M in pre-seed funding} to: build the team (\$300K--\$500K),
complete atlas reprocessing (\$21K--\$35K), train models (\$52K--\$143K), LLM
annotation (\$20K--\$50K), launch singlet.ai (\$67K--\$163K), and develop
clinical PGI (\$60K--\$125K).

\paragraph{Return Thesis.} Stage~1 (Year~1--3): \$4.5M~ARR, \$45M valuation,
15$\times$ on pre-seed. Stage~2 (Year~3--5): \$19M~ARR, \$190M valuation,
60$\times$ on pre-seed.

%% ──────────────────────────────────────────────────────────────────
\section{Customer Segments}
\label{sec:customers}

\subsection{Academic Researchers}

Academic and research laboratories represent 70.8\% of the \$4.89B single-cell
analysis market. Key pain points include: data acquisition requiring weeks of
format wrangling across 48+ types, batch effects dominating cross-dataset
analyses, inconsistent cell type annotation, compute requirements exceeding lab
budgets, reproducibility crisis from dependency rot, and foundation model
confusion (scGPT, Geneformer, UCE---none interoperable).

Value proposition: \texttt{singlet.query(species="human", tissue="lung",
disease="IPF")} returns uniformly processed, annotated AnnData, eliminating the
2--6~week download and harmonization cycle. NMF gene programs provide
interpretable biological discovery. Cross-species comparison across 200+ species
is unmatched.

\subsection{Biotech and AI-Native Companies}

The AI drug discovery market is projected at \$13.8B by 2030 (47\% CAGR). Every
foundation model team needs the same thing: massive, clean, uniformly processed
training data. PyTorch DataLoaders that stream directly from compressed matrices
are essential---current workflows die at scale. NMF pre-embeddings (50--100
features vs.\ 20,000 genes) offer 200$\times$ faster training.

\subsection{Core Facilities}

$\sim$500 genomics core facilities in the US alone, each supporting 20--100 PIs.
Analysis bottleneck dominates: sequencing turnaround is 1--3 days, analysis
turnaround is 2--6~weeks. Atlas-referenced QC scoring
(\texttt{singlet.benchmark\_qc()}) provides objective, defensible quality
assessment at percentile resolution against hundreds of comparable datasets.

\subsection{Pharmaceutical Companies}

Pharma routinely pays \$50K--\$500K/yr for data products. Internal data
harmonization projects cost \$500K+ in FTE time. Cross-species translation
(mouse~$\to$ human) is critical and unsolved---our 200+ species NMF factor space
directly addresses this. Target Explorer API
(\texttt{singlet.query\_gene("TREM2")}) returns cell-type-resolved expression
across all published studies, replacing 2--4~weeks of FTE time per gene.

%% ──────────────────────────────────────────────────────────────────
\section{Conclusion}
\label{sec:conclusion}

Singlet AI occupies a unique position at the intersection of four convergent
advances: sequence-to-function genomic models, paired WGS + single-cell data,
atlas-scale uniformly reprocessed single-cell repositories, and biological
reasoning LLMs. Our six-layer defense---from the open-source sequence grammar
trunk through our proprietary single-cell prediction head, multimodal atlas,
interpretable NMF models, reasoning AI, and clinical product---creates a
defensible path from pre-seed to clinical deployment.

The asymmetric bet: pre-seed funds the data platform (the proven Tempus-like
model). The clinical PGI upside---AI agents that explain any genome to any
clinician---is optionality on top.

\bibliographystyle{unsrtnat}
\bibliography{refs}

\end{document}
